% !PRINTED 227
\weberpage{227}
%<<CONT
If $m, n$ are any two coprime numbers, the latter odd, and if $\mu$ further denotes the number of those among the numbers

\[
m,\ 2m,\ 3m \cdots \frac{n-1}{2}\,m,
\]

whose absolutely least residues $(\mathrm{mod}\ n)$ are negative, then

\weq{8}{(-1)^\mu = \left(\frac{m}{n}\right),}

where $\left(\dfrac{m}{n}\right)$ is the \gesperrt{Legendre-Jacobi} symbol from the theory of quadratic residues\textsuperscript{1)}.

This theorem now leads immediately to the determination of the sign in (7). For if the function $\Phi(x)$ in (5) is odd, then in the passage from $\Omega$ to $m\,\Omega$ exactly $\mu$ factors in (5) change their sign, and we conclude

\weq{9}{\Pi(m\,\Omega) = \left(\frac{m}{n}\right) \Pi(\Omega).}

If $n$ is not a square number, then $m$ can always be assumed such that $\left(\dfrac{m}{n}\right) = -1$. For if $n = p^\lambda n'$, with $\lambda$ odd, $n'$ not divisible by $p$, and $\beta$ a quadratic non-residue of $p$, then one need only determine $m$ from the congruences

\[
m \equiv \beta \pmod{p^\lambda}, \quad m \equiv 1 \pmod{n'}
\]

in order to find such a number $m$. And then $\Pi(m\,\Omega) = -\Pi(\Omega)$, and not $\Pi(\Omega)$, but only $\Pi(\Omega)^2$ is a root of a transformation equation. But if $n$ is a square number, then $\Pi(m\,\Omega) = \Pi(\Omega)$ for every $m$, and consequently $\Pi(\Omega)$ itself is a root of a transformation equation. Hence the theorem follows:

\begin{quote}
\gesperrt{For an odd function $\Phi$ it is $\Pi(\Omega)^2$, and only when $n$ is a square number $\Pi(\Omega)$ itself, that is a root of a transformation equation\textsuperscript{2)}.}
\end{quote}

\textsuperscript{1)}~Cf.\ on this theorem: \gesperrt{Schering} and \gesperrt{Kronecker} in the Monatsbericht der Berliner Akademie of 22 June 1876. \gesperrt{Schering}, Acta mathematica I.

\textsuperscript{2)}~This simplification of the multiplier equation in the case where $n$ is a square was discovered by \gesperrt{Joubert}, but was proved by a method entirely different from ours. ``Sur les équations, qui se rencontrent dans la théorie de la transformation des fonctions elliptiques.'' Paris 1876.
